\begin{agradecimentos}
À Universidade Federal de Sergipe (UFS), em especial ao Centro de Ciências Exatas e Tecnologia (CCET) e ao Departamento de Computação (DCOMP), pela sólida formação científica, tecnológica e ética proporcionada ao longo de toda a graduação em Engenharia de Computação.

Ao meu orientador, Prof. Dr. Gilton José Ferreira da Silva, pela confiança depositada, pelas diretrizes acadêmicas precisas, pelo incentivo contínuo ao rigor metodológico e pela visão integradora entre a computação e a sociedade.

À minha coorientadora, Prof.ª Dr.ª Ana Waleska Silva Patrício, do Departamento de Enfermagem da UFS, pela generosidade no compartilhamento de seu vasto conhecimento clínico em segurança do paciente e pela liderança inspiradora que possibilitou a perfeita convergência interdisciplinar deste projeto.

Aos docentes e colaboradores do DCOMP e do Departamento de Enfermagem, cujos ensinamentos e questionamentos enriqueceram as bases teóricas e práticas deste trabalho.

À equipe multiprofissional do Hospital Universitário da UFS (HU-UFS/EBSERH), que diariamente vivencia os desafios da assistência e inspira o desenvolvimento de soluções tecnológicas de impacto direto na qualidade do cuidado.

Aos meus pais e familiares, pelo amor constante, pela paciência nas noites de estudo e desenvolvimento e por serem a fonte inesgotável de motivação para todos os meus sonhos e conquistas.

Aos colegas de curso e amigos da UFS, pelos debates enriquecedores, pelo companheirismo nos laboratórios e pela amizade fraterna cultivada durante esta marcante trajetória acadêmica.
\end{agradecimentos}